% =====================================================================
% Section 7 — Automated AI Workflow: Methods (REWRITE)
% Table-referencing, pattern-focused. Paste-ready LaTeX prose.
%
% Table label conventions used below (rename \ref keys to match your .tex):
%   Table 22 -> \label{tab:main-audit-scale}
%   Table 23 -> \label{tab:main-resolver-validation}
%   Table 24 -> \label{tab:main-breakdowns-across-levels}
%
% NUMBERS TO RE-VERIFY against the freshly regenerated tables before submit
% are marked with %%VERIFY. Everything else is a stable headline figure.
% =====================================================================

\section{Automated AI Workflow: Methods}

We evaluated the automated audit workflow at four increasingly larger scales:
a rerun of the original 80-record Phase~I sample, a rerun of the larger 2025 DAP
validation set, an uncontrolled run over the full Ceska / Observatory Hill record
set, and an uncontrolled run over the entire museum fungal collection. These
levels trade evidence quality for scale. The two smaller levels have a
human-reviewed ground truth and let us measure resolver accuracy; the two larger
levels have none and instead estimate the maintenance burden a curator would face
across a whole collection. Table~\ref{tab:main-audit-scale} summarizes the sample
frame, purpose, and principal limitation of each level;
Table~\ref{tab:main-resolver-validation} reports resolver performance wherever a
ground truth existed; and Table~\ref{tab:main-breakdowns-across-levels} reports
how the seven-category breakdown typology distributes at each scale.

Read together, the three tables tell a single story. Explicit cross-references
between the museum and Mushroom Observer are rare and, where they exist, almost
never mutual; the resolver recovers most human-reviewed links but conservatively
defers exactly the cases the manual audit also flagged for physical or curatorial
judgment; and the relative weight of the seven breakdown categories shifts
systematically as the audit widens from a hand-picked sample to the full
collection. We take each level in turn and then return to the cross-scale pattern.

\subsection{Recreating the Phase~I sample}

We reran the automated audit on the exact 80 Mushroom Observer records used in the
manual Phase~I sample. Of these, 34 had at least one manually listed museum
catalog candidate and 46 had no manually identified museum counterpart. Against
the 34 manual-candidate rows the resolver achieved 73.5\% recall and 100.0\%
precision among the links it selected (Table~\ref{tab:main-resolver-validation}).

The Phase~I rerun shows why the automated audit should be read as a diagnostic
instrument rather than a replacement for manual harmonization. The records it
missed or deferred were precisely the kinds the manual workflow set aside for
physical inspection, image comparison, or curator judgment: uncertain catalog
numbers, multiple candidate museum specimens, duplicated Mushroom Observer
records, and date or taxonomic inconsistencies. Among the 46 records without a
manually listed candidate, the resolver proposed three possible museum matches.
We treat these not as confirmed harmonizations but as a review queue that can help
curators prioritize unresolved records.

\subsection{Recreating the DAP 2025 audit}

Because the DAP dataset begins with Mushroom Observer records, this run used a
top-down audit direction: each Mushroom Observer record was traced outward to the
museum database, MyCoPortal, GBIF, and GenBank where possible. The DAP file
contained 408 Mushroom Observer records, of which 355 carried a DAP-listed museum
catalog number. Assessing the current exported records for \emph{explicit}
cross-references reproduces the paper's central linkage finding at larger scale:
no DAP pair was bidirectionally cross-referenced, and where a link existed at all
it ran overwhelmingly in one direction, from the Mushroom Observer record to the
museum, with the reciprocal citation on the museum record absent
(Table~\ref{tab:main-breakdowns-across-levels}, HB1). Across the 408 records the
audit surfaced all seven breakdown types.

The DAP set's value is that it also supplies a ground truth for what \emph{should}
be linked, letting us validate the resolver's matching component rather than only
its cross-reference bookkeeping. For the paper-facing validation we excluded 33
rows in which the raw DAP export repeats a single identifier across many taxa and
dates, leaving 322 validation links. On these the resolver selected the same
museum catalog number as DAP for 277 records, a different number for 5, and no
acceptable match for 40, giving 86.0\% recall and 98.2\% precision
(Table~\ref{tab:main-resolver-validation}). Critically, the 40 misses were not
absent from the museum data; they failed the resolver's matching rules, splitting
into names below the similarity threshold despite matching date and locality,
genus- or name-blocked pairs, and matching names with conflicting collection
dates.
The resolver additionally proposed museum candidates for 20 of the 53 Mushroom
Observer records that had no DAP-listed catalog number. These failure modes are
instructive: date conflicts, where different collection conventions place the same
specimen on different dates, and name drift, where taxonomy has moved since
collection, are exactly the reconciliations that depend on curatorial knowledge or
an up-to-date synonym source rather than string matching.

\subsection{Audit of C. / O.H. records}

Using the same skill-supported framework, we ran the audit over the broader
Ceska / Observatory Hill record set, reported as two frames in
Table~\ref{tab:main-audit-scale}: a museum side scoped by locality (Observatory
Hill) and a Mushroom Observer side scoped by both locality and collector. The two
sides are scoped differently because, on the museum export, neither the collector
nor the locality field supports a clean programmatic filter. Collector-based
scoping is unreliable in two compounding ways. First, many of C.'s specimens were
collected jointly and are recorded with co-collectors---for example ``Oluna Ceska;
Paul Kroeger'' or attributed to a collecting group such as SVIMS---so a search on
her name alone silently omits every specimen on which she is one of several
collectors, skewing the count of Observatory Hill holdings downward. Second, in the
current export the structured collector field is not reliably populated, so it
cannot be used to filter programmatically at all. Locality-based scoping is also
only approximate: Observatory Hill is not entered as a single locality value but as
a set of descriptive sub-localities naming different sections of the site, so a
locality filter captures the area only as well as those free-text descriptions
allow. We therefore scope the museum side by locality and treat the resulting frame
as an approximation of C.'s Observatory Hill holdings rather than an exact
enumeration. This asymmetry reflects how the source records were entered, not an
analytical choice, and the museum-side counts should be read with it in mind.

At this scale the cross-reference pattern from the smaller levels holds and
sharpens: only a handful of Observatory Hill museum records carry a Mushroom
Observer identifier, yet hundreds are cited by at least one Mushroom Observer
record, so the linkage is almost entirely one-directional from the citizen-science
side to the museum (Table~\ref{tab:main-breakdowns-across-levels}). The downstream
harvested platforms behave very differently: MyCoPortal coverage of the
Observatory Hill museum records is near-complete, GBIF coverage is weaker, and
GenBank remains barely integrated, with most associated accessions uncited on the
museum record. The resolver proposed at least one museum candidate for roughly a
third of the Mushroom Observer records and left the majority with no digital
candidate (Table~\ref{tab:main-resolver-validation}). The system can therefore
scale explicit cross-reference checks, flag candidate matches, and produce review
queues across thousands of records, but the large residue of unmatched and
ambiguous cases is itself the finding: visual, physical, and curatorial judgment
remain necessary for high-confidence harmonization.

\subsection{Full museum fungal collection run}

The full-collection run covers all 34,856 museum fungal records and is a
collection-scale triage rather than an accuracy test, since no collection-wide
manual gold standard exists. Its headline is the scarcity of mutual linkage: only
17 records in the entire collection are bidirectionally cross-referenced with
Mushroom Observer, and although the museum's harvested links to MyCoPortal and
GBIF are largely intact, its connection to derivative sequence data is not---within
the locally available GenBank evidence, 212 of 213 accessions associated with a
museum catalog number are not cited on the corresponding museum record. The
resolver generated 2,444 metadata-based museum--Mushroom Observer candidate pairs,
which we report as a review queue rather than confirmed identities. Mapped onto the
seven-category typology (Table~\ref{tab:main-breakdowns-across-levels}), the
full-collection audit flags records under every breakdown, dominated by
missing-cross-reference, repository-absence, and duplicate-candidate cases.

\subsection{Patterns across scales}

Two movements stand out when the breakdown proportions are compared across scales
(Table~\ref{tab:main-breakdowns-across-levels}). First, absence from repositories
(HB3) falls sharply as the sample widens, from 58.8\% in the Phase~I automated
rerun to 5.6\% collection-wide. This suggests the high backlog rate we measured by
hand is a property of the C. collection's particular digitization history rather
than a collection-wide condition: the full export is, on the whole, better
digitized than the sample we deliberately drew from the backlog. Second, duplicate
record candidates (HB6) move in the opposite direction, from single digits in the
manual audit to 28.2\% collection-wide. This is more consistent with an artifact of
automated harvesting and resolver behavior at scale---near-duplicates the resolver
cannot yet disambiguate---than with a real rise in true duplication. Whether HB6 at
scale is a data-quality problem or a resolver-precision problem changes what should
be done about it, and future work should test the two apart. We treat both
movements as hypotheses the resolver's output makes newly visible, not as confirmed
findings, because neither larger level had a curator-reviewed ground truth to check
them against.
